
\section{Introducción al estado del arte}

El problema de este trabajo es que un usuario no puede saber a dónde lo lleva una URL
antes de que su teléfono la abra al escanear un código QR. Ese problema cruza cuatro
puntos de investigación que hasta ahora han avanzado bastante por separado: (i) el
estudio del código QR como vector de ataque y qué tan susceptibles son los usuarios;
(ii) la detección automática de URLs maliciosas, desde listas negras hasta
clasificadores que analizan cómo está escrita la URL; (iii) los sistemas de
reputación y las plataformas colaborativas de \emph{threat intelligence}; y (iv) el
campo de \emph{usable security}, que estudia cuándo una advertencia de seguridad
realmente cambia lo que hace la persona que la recibe.

Al revisarlas juntas aparece un desbalance que explica por qué este proyecto tiene
sentido. Los tres primeros puntos están bien desarrollados y con resultados medibles:
hay taxonomías completas de ataques sobre QR \cite{krombholz2014qr}, estudios de
campo sobre susceptibilidad \cite{vidas2013qrishing}, clasificadores que reportan
precisiones arriba del 95~\% \cite{ma2009beyond,le2011phishdef} y una infraestructura
de reputación de URLs funcionando a escala mundial
\cite{google_safebrowsing,virustotal_apiv3,phishtank_devinfo,urlhaus_api,misp_project}.
El cuarto punto, en cambio, ha mostrado una y otra vez que detectar bien no basta para
proteger: entre el 25~\% y el 70~\% de los usuarios ignora la advertencia y sigue
adelante, según el tipo de amenaza \cite{akhawe2013alice}.

Lo que esto implica es que hoy el \emph{quishing} no está limitado por la falta de
capacidad para detectar, sino porque esa capacidad no se aplica en el momento y en el
aparato donde ocurre el ataque, y porque el resultado no se comunica de forma que el
usuario pueda hacer algo con él. Este capítulo revisa los trabajos de cada punto,
compara sus enfoques y resultados, señala qué aporta y qué le falta a cada uno, y
delimita el vacío que el prototipo busca llenar.

\subsection*{Procedimiento de la revisión}

Varias conclusiones de este capítulo dicen que algo no existe en la literatura. Para
que esas afirmaciones se puedan verificar y quede claro hasta dónde llegan, conviene
explicar cómo se hizo la búsqueda.

Se buscó en IEEE~Xplore, ACM Digital Library, Scopus, SpringerLink, arXiv y Google
Scholar, además de la documentación oficial de los proveedores de reputación y de las
plataformas móviles, y de informes públicos de la industria de seguridad. Los términos
de búsqueda combinaron 	exttt{"QR code" AND (phishing OR quishing)},
	exttt{quishing}, 	exttt{"malicious URL detection" AND (mobile OR lexical)},
	exttt{"phishing blacklist" AND (coverage OR "zero-hour")},
	exttt{"security warning" AND (effectiveness OR usable security) AND mobile} y
	exttt{"URL reputation" AND (aggregation OR "multi-source")}, en inglés y en español.
El rango de años fue 2008--2025, ampliado hacia atrás solo para incluir trabajos
fundacionales que siguen vigentes por la cantidad de veces que se citan.

Se incluyeron trabajos publicados en conferencias o revistas indexadas, estándares
internacionales y documentación oficial de las tecnologías que se usan en el
prototipo. Se dejaron fuera los trabajos sobre generación o procesamiento de imagen
del símbolo QR sin componente de seguridad, las propuestas de firmar códigos QR con
esquemas criptográficos que obligan a cambiar el estándar o la infraestructura de
quien emite el código, y las publicaciones sin acceso al texto completo. La selección
se hizo leyendo título y resumen, y después el texto completo de los que quedaron.

Esto no es una revisión sistemática formal: no hubo doble revisión independiente ni
protocolo registrado de antemano. Por eso, cuando más adelante se diga que algo no se
encontró, hay que leerlo con esa precisión: no se afirma que esa evidencia no exista,
sino que no apareció con la búsqueda descrita aquí. Es una limitación de esta revisión
y se declara como tal.

Hay un segundo punto que conviene aclarar desde el inicio. Buena parte de lo que se
sabe sobre \emph{quishing}, sobre todo en cuanto a volúmenes y crecimiento, viene de
informes de empresas de seguridad \cite{verizon2024dbir,insikt2024qrcode,unit42_2023qr}
y no de publicaciones arbitradas. Esas fuentes sirven porque tienen acceso a
telemetría que nadie más tiene, pero no pasan por revisión de pares y casi nunca
explican su metodología con suficiente detalle como para reproducirla. Por eso en este
capítulo se usan para dimensionar el fenómeno, no para sostener decisiones técnicas;
esas se apoyan en publicaciones indexadas, estándares y documentación oficial.

\section{Revisión de antecedentes relevantes}

\subsection{El código QR como vector de ataque: caracterización y evidencia empírica}

El trabajo que abrió este punto es el de Krombholz et al. \cite{krombholz2014qr},
la primera taxonomía ordenada de ataques sobre códigos QR vista desde
\emph{usable security}. Su aporte principal fue separar dos cosas: los ataques que
manipulan el código físico (pegarle una calcomanía encima, reimprimir un cartel) y los
que aprovechan lo que está codificado adentro. Identificaron las URLs como el
contenido más peligroso y más común y, lo más importante para este trabajo, plantearon
el problema como un asunto de usabilidad y no solo de criptografía o de integridad del
código: el usuario no puede verificar lo que no puede leer.

La evidencia empírica clave viene de Vidas et al. \cite{vidas2013qrishing}. En su
estudio \emph{QRishing} pegaron códigos QR en lugares públicos reales y midieron
cuánta gente los escaneaba de verdad. Encontraron que los usuarios de teléfonos
inteligentes caen bastante más en phishing por QR que en phishing por enlace directo,
lo que confirmó con números lo que Krombholz et al. habían planteado a nivel
conceptual. Metodológicamente destaca porque midieron conducta real en la calle en
lugar de preguntarle a la gente qué haría, evitando así que las respuestas salieran
sesgadas por querer quedar bien. También documentaron que la mayoría escaneó por pura
curiosidad, sin fijarse antes en el destino.

Los dos trabajos comparten un problema de fecha: son de 2013 y 2014, o sea antes de
que la pandemia de COVID-19 disparara el uso de códigos QR a partir de 2020
\cite{salahdine2019social}. En el escenario que ellos describen, escanear un QR
requería bajar una app aparte; hoy la cámara del teléfono lo hace sola, lo que quita
todavía más fricción y hace menos probable que alguien se detenga a pensar. Sus
conclusiones sobre susceptibilidad no solo siguen siendo válidas, sino que
probablemente el problema empeoró; pero las cifras concretas que reportan hay que
tomarlas como piso, no como estimación actual.

Para saber cómo está el fenómeno hoy hay que recurrir a informes de la industria.
Insikt Group \cite{insikt2024qrcode} registró un aumento cercano al 248~\% en menciones
a \emph{quishing} en fuentes de inteligencia entre el tercer y el cuarto trimestre de
2023 y, más importante que la cifra, documentó que el ataque se integró a plataformas
de \emph{phishing-as-a-service} como Tycoon 2FA y Greatness. Eso significa que ya no
hace falta saber hacerlo: cualquiera puede comprar el servicio, lo que amplía muchísimo
el número de atacantes posibles. Unit~42 de Palo Alto Networks \cite{unit42_2023qr}
explicó el incentivo: plataformas como Microsoft 365 y Google Workspace revisan el
texto de los correos pero no decodifican los códigos QR que vienen como imagen, así
que el QR funciona como forma de esquivar esos filtros.

Lo valioso de estos informes no son tanto sus cifras de crecimiento, cuya
metodología no se puede reproducir y cuyas bases de comparación cambian de un reporte a
otro, sino cómo describen el incentivo económico y la forma de evadir los filtros. Eso
sí se puede verificar por otro lado, y es lo que sostiene la premisa central de este
trabajo: si el filtro de correo no puede leer el QR, el único lugar donde la URL se
vuelve legible es el teléfono del usuario al momento de escanear.

\subsection{Detección de URLs maliciosas: listas negras, heurísticas y aprendizaje automático}

La referencia obligada aquí es Sahoo, Liu y Hoi \cite{sahoo2017malicious}, que
ordenan las técnicas de detección en tres grupos (listas negras, heurísticas y
aprendizaje automático) y explican qué gana y qué pierde cada una en cobertura,
latencia, falsos positivos y capacidad de detectar amenazas nuevas. Su aporte es sobre
todo organizativo y sigue siendo el marco más usado; su límite es la fecha, anterior a
que se generalizaran los transformadores para clasificar texto, así que lo que dicen
sobre aprendizaje profundo quedó desactualizado.

El trabajo de Sheng et al. \cite{sheng2009empirical} es la evaluación más citada
sobre los límites de las listas negras. Midieron cuánto cubrían varias listas en
operación y cuánto tardaban en incorporar sitios nuevos, y encontraron que las más
grandes alcanzaban alrededor del 20~\% de los sitios de phishing en el momento en que
la URL se creaba, con una cobertura que mejoraba conforme el sitio llevaba más tiempo
circulando y siendo denunciado. De ahí salió el concepto de \emph{zero-hour phishing}
como límite estructural de este enfoque.

Ahora bien, aplicar ese hallazgo tal cual al día de hoy sería un error, y vale la
pena señalarlo porque afecta directamente cómo se justifica el filtro F1 de este
proyecto. El estudio es de 2009, cuando los sistemas de reputación eran mucho menos
maduros; las plataformas actuales actualizan sus bases en cuestión de minutos con
análisis automatizado a gran escala. El 20~\% no se puede trasladar al presente. Lo que
sí sigue vigente, y es lo que sostiene el argumento, es la limitación de fondo: los
informes de APWG de 2024 y 2025 \cite{apwg2024q1,apwg2025q1} muestran que los sitios
maliciosos que duran menos de una hora antes de ser indexados siguen siendo una parte
importante de los ataques, con más de un millón de ataques únicos solo en el primer
trimestre de 2024. Ninguna lista puede contener algo que todavía no ha sido visto, por
más rápido que se actualice. La conclusión correcta no es que las listas negras ya no
sirvan, sino que son necesarias pero no suficientes: cubren muy bien y muy rápido lo
que ya se conoce, y hay que complementarlas con algo que evalúe lo que no.

El tamaño de esa brecha se puede estimar con los mismos informes de APWG
\cite{apwg2024q1,apwg2025q1}, que documentan un flujo constante de dominios de
phishing recién creados y una porción importante de sitios que duran menos de lo que
toma indexarlos. Mientras más ataques usen dominios registrados especialmente para una
campaña, menos alcanza a cubrir cualquier lista en el instante en que el usuario
escanea.

Es probable que en el \emph{quishing} esa proporción sea al menos parecida a la del
phishing por enlace directo, porque repartir códigos QR físicamente favorece campañas
cortas y localizadas, para las que registrar un dominio propio resulta cómodo. Hay que
subrayar, eso sí, que esto es un supuesto derivado de cómo opera el ataque y no una
medición: la revisión no encontró estudios que cuantifiquen específicamente cuántos
dominios únicos se usan en campañas de \emph{quishing}. Se declara como supuesto y no
se usa para sostener ninguna decisión de diseño; solo indica hacia dónde apunta el
efecto esperado.

Frente a este límite, el enfoque que clasifica URLs por sus características es el más
prometedora para un contexto móvil. Ma et al. \cite{ma2009beyond}, en un trabajo de
KDD~'09 que abrió el camino, mostraron que con datos sacados del nombre de dominio y
de cómo está escrita la URL se puede llegar a precisiones cercanas al 99~\% detectando
phishing conocido, sin necesidad de visitar el sitio. Le, Markopoulou y Faloutsos
\cite{le2011phishdef} llevaron la idea al extremo con \emph{PhishDef}, que usa
únicamente características de la URL y reporta precisiones arriba del 95~\%, con la
ventaja de degradarse de forma controlada cuando los datos vienen con ruido, gracias al
aprendizaje en línea.

Lo que hace decisivos estos trabajos para una app móvil no es su precisión, sino algo
que se desprende de no tener que abrir el sitio. De ahí salen tres ventajas. No hay que
esperar una petición HTTP al destino, que además puede ser malicioso. No se expone la
infraestructura del analizador ante ese destino. Y, sobre todo, no se le avisa a nadie
más que un usuario específico está por entrar a una URL específica. Esto último es un
tema de privacidad que la literatura revisada apenas menciona de pasada, pero que en un
producto para usuarios finales pesa bastante más.

Los límites de este enfoque también importan. Los resultados salieron de conjuntos de
datos históricos y en condiciones de laboratorio; ninguno de los trabajos revisados
reporta cómo se comporta en uso real sobre un teléfono, ni bajo las restricciones de
cómputo y batería que eso implica. Además, las características de la URL pierden
efectividad ante técnicas de evasión que trabajos posteriores ya documentaron: usar
acortadores, aprovechar dominios legítimos comprometidos o generar nombres por
algoritmo que estadísticamente se ven igual de inocentes que los normales. Por eso este
trabajo pone el análisis heurístico en un filtro posterior (F2) y no como único
mecanismo.

\subsection{Sistemas de reputación y plataformas colaborativas de threat intelligence}

La infraestructura en la que se apoya el filtro F1 está documentada sobre todo en
manuales técnicos oficiales y no en publicaciones arbitradas, lo cual ya dice algo
sobre el estado del arte en esta área.

Google Safe Browsing \cite{google_safebrowsing} es el sistema de reputación con más
URLs indexadas, y es el que alimenta las advertencias de Chrome, Safari y Firefox. Su
Lookup API~v5 permite consultar URLs una por una en tiempo real, normalmente en menos
de 200~ms. Lo que aporta al diseño propuesto es llevar al escaneo móvil una cobertura
que ya existe en el navegador de escritorio; su desventaja es que no se sabe con qué
criterio clasifica, así que no hay forma de auditar por qué marcó o no marcó una
URL.

VirusTotal \cite{virustotal_apiv3} hace algo distinto y complementario: en vez de
mantener su propia lista, junta los veredictos de más de setenta motores antivirus y
escáneres de URL. La ventaja es que se reducen los falsos negativos, porque una URL que
se le escapa a un motor puede ser detectada por otro. La desventaja pesa mucho en un
prototipo académico: el plan gratuito permite solo cuatro consultas por minuto y
quinientas al día. Ese tope limita cuánta validación empírica se puede hacer y obliga a
diseñar el filtro para que siga funcionando cuando esta fuente no esté
disponible.

PhishTank \cite{phishtank_devinfo}, de Cisco Talos, y URLhaus \cite{urlhaus_api}, de
abuse.ch, funcionan de forma colaborativa: la comunidad reporta las URLs y estas se
verifican por votación o curaduría. Lo que aportan es rapidez para incorporar campañas
nuevas, más que los sistemas totalmente automatizados en ciertos nichos. Su límite es
que dependen de qué tan activa esté la comunidad, lo que hace que la cobertura sea
despareja según la región y el idioma. Eso pega directo en el contexto guatemalteco y
latinoamericano de este trabajo, donde seguramente se reporta bastante menos que en
inglés. MISP \cite{misp_project} completa el panorama como el estándar de hecho para
intercambiar indicadores de compromiso. Su relevancia aquí es indirecta pero real:
varias de las fuentes integradas exportan o consumen datos en formatos compatibles con
MISP, y eso explica que el ecosistema funcione junto.

Al mirar estas plataformas en conjunto aparece un vacío que ninguna documentación
oficial cubre: en la literatura revisada no hay un estudio independiente que compare
qué tanto coinciden estas fuentes sobre un mismo conjunto de URLs, ni que describa cómo
se comporta un veredicto combinado frente a los veredictos por separado. Cada proveedor
publica lo que él hace; nadie publica cómo funcionan juntos. La validación del filtro F1
prevista en este trabajo, aunque sea de alcance modesto, aporta algo en esa
dirección.

\subsection{Soluciones comerciales de escaneo seguro de QR}

En las tiendas de apps ya hay soluciones que atacan parte del problema. Herramientas
como \emph{Safe QR Scan -- Secure QR Scanner} y \emph{SafeQR Reader $\cdot$ Block Scams},
disponibles en la App Store, revisan la URL antes de redirigir, comparándola contra
bases de datos de dominios maliciosos conocidos.

Su mérito es haber demostrado en el mercado que el punto de intervención funciona:
confirman que meterse entre el escaneo y la apertura del navegador es viable
técnicamente y que a los usuarios no les molesta. Pero tienen tres limitaciones de
fondo. Primero, como ya documentó la literatura revisada
\cite{sheng2009empirical,vidas2013qrishing}, trabajan casi todas con listas negras, que
no alcanzan contra dominios recién creados, justamente los que dominan en campañas de
\emph{quishing}. Segundo, son productos cerrados: no publican su arquitectura, ni sus
fuentes de datos, ni sus tasas de error, así que no hay manera de evaluarlas de forma
independiente ni de compararlas científicamente. Tercero, ninguna explica en qué basó
sus decisiones de interfaz, aun cuando la literatura de \emph{usable security} dice que
ahí es donde fallan la mayoría de sistemas de advertencia.

Conviene aclarar que este trabajo no incluye una comparación formal contra estas
apps comerciales; esa exclusión está declarada en el capítulo de alcance. Se revisan
aquí para mostrar cómo está la práctica actual y para dejar claro que lo que se propone
--una arquitectura documentada, con fuentes declaradas y decisiones de interfaz
fundamentadas-- juega en otro plano que el de competir en funciones con productos de
mercado.

\subsection{Usable security y efectividad de advertencias de seguridad}

Este punto aporta el fundamento de las decisiones de interfaz. Puede sostenerse,
además, que es la que explica por qué la sola disponibilidad de detección eficaz no ha
resuelto el problema.

Cranor \cite{cranor2008framework} propuso el marco del humano en el ciclo, que
descompone el comportamiento seguro en tres condiciones encadenadas: la persona tiene
que notar la señal, entender qué significa y actuar en consecuencia. Sirve como
diagnóstico, porque permite ubicar en cuál de los tres eslabones está fallando un
sistema concreto en lugar de echarle la culpa genéricamente al descuido del usuario. Su
límite es que es un modelo conceptual y no un instrumento de medición, así que
validarlo directamente es complicado.

Egelman, Cranor y Hong \cite{egelman2008warned} aportaron evidencia experimental
sobre qué hace efectiva una advertencia de phishing en el navegador: interrumpir lo que
la persona está haciendo, destacar visualmente, usar lenguaje concreto en vez de
técnico y aparecer en el momento en que más importa. Akhawe y Felt
\cite{akhawe2013alice} llevaron el análisis a datos reales de Chrome, con un hallazgo
muy relevante para este proyecto: la tasa de continuación, o sea cuánta gente ignora la
advertencia y sigue de todos modos, fue de alrededor del 70~\% en advertencias de SSL
frente a cerca del 25~\% en advertencias de malware. Esa diferencia muestra que la
reacción depende de qué tan grave se percibe la amenaza y de qué tan bien se entiende,
no de que la gente ignore los avisos por costumbre. La implicación de diseño es
directa: distinguir claramente entre niveles de veredicto no es un detalle estético,
es un requisito para que la advertencia funcione.

El problema común de este punto es el contexto: los tres trabajos estudian
advertencias en navegadores de escritorio. Nadie ha validado si sus hallazgos se
trasladan al entorno móvil, que es distinto: pantalla chica, interacción táctil, uso
frecuente en movimiento y con prisa, y sin barra de direcciones para verificar. Es,
entonces, una extrapolación razonable pero no comprobada, y este trabajo la asume
explícitamente como tal.

\subsection{Estándares y documentación técnica de referencia}

El estándar ISO/IEC~18004 \cite{iso18004_2015,iso18004_2024} define cómo funciona el
código QR: modos de codificación, niveles de corrección de errores y versiones de
tamaño. Importa como antecedente por dos razones. Por un lado, delimita técnicamente
qué puede llevar adentro un QR y, por lo tanto, qué tiene que clasificar el módulo de
escaneo antes de decidir si corresponde analizarlo. Por otro, y esto pesa más para el
planteamiento del problema, el estándar no incluye ningún mecanismo para autenticar
quién generó el código ni para verificar que el contenido no fue alterado. No es un
defecto de la norma, porque se diseñó en 1994 para trazabilidad industrial
\cite{iso18004_2015} y no como canal de confianza, pero ahí está la raíz técnica de la
vulnerabilidad: dentro del estándar no hay forma de distinguir un QR legítimo de uno
pegado encima. Cualquier solución tiene que ubicarse por encima de la capa del código,
lo que respalda la decisión de este trabajo de analizar la URL en texto (RFC~3986) y no
el símbolo QR.

El OWASP MASVS~v2.1 \cite{owasp2024masvs} da el marco para verificar la seguridad de
la app misma, con categorías que aplican directo al prototipo: MASVS-STORAGE para el
historial local cifrado, MASVS-NETWORK para la comunicación con el backend y las APIs, y
MASVS-PLATFORM para el manejo de permisos de cámara y ubicación. Lo útil es que ofrece
criterios verificables en lugar de buenas prácticas sueltas; su límite reconocido es que
es un estándar de requisitos y no una guía de implementación. La documentación oficial de
cada plataforma (App Transport Security \cite{apple_ats} y Apple Platform Security
\cite{apple2024platformsecurity} en iOS, Network Security Configuration
\cite{android_netsecconfig} en Android) complementa el estándar con los mecanismos
concretos.

Por último, la documentación de las librerías y APIs nativas de escaneo
(VNDetectBarcodesRequest del framework Vision \cite{apple_vndetectbarcodes}, ML~Kit
Barcode Scanning \cite{google_mlkit_barcode} y react-native-vision-camera
\cite{rousavy_visioncamera}) confirma que se puede procesar todo en el teléfono en menos
de 100~ms. Este dato es el que hace posible el objetivo de mantener la latencia total
por debajo de cinco segundos: si hubiera que mandar la imagen a un servidor para
decodificarla, ni el tiempo ni la privacidad serían alcanzables.

\section{Comparación de los trabajos revisados}

Poniendo los antecedentes uno junto a otro, se pueden ordenar según en qué momento
del ataque intervienen y qué tipo de evidencia aportan.

	extbf{Por enfoque metodológico} se distinguen cuatro grupos con características
distintas. Los estudios de campo que miden conducta real \cite{vidas2013qrishing} son
los más cercanos a la realidad, pero con muestras pequeñas y difíciles de repetir. Las
evaluaciones sobre conjuntos de datos
\cite{sheng2009empirical,ma2009beyond,le2011phishdef} dan métricas precisas y
comparables, aunque en condiciones de laboratorio que no reflejan el uso real. Los
estudios a gran escala con datos de producción \cite{akhawe2013alice} combinan volumen y
realismo, pero solo los puede hacer quien controla la plataforma. Y los informes de la
industria \cite{insikt2024qrcode,unit42_2023qr} ofrecen la foto más actual del panorama
de amenazas, a cambio de no pasar por revisión de pares ni ser reproducibles.

	extbf{Por punto de intervención} la comparación dice más. Los sistemas de reputación
\cite{google_safebrowsing,virustotal_apiv3,phishtank_devinfo,urlhaus_api} actúan sobre la
URL una vez que el ecosistema ya la conoce; cubren mucho y responden rápido, pero no ven
lo nuevo. Los clasificadores léxicos \cite{ma2009beyond,le2011phishdef} actúan sin
importar el historial de la URL, así que pueden evaluar dominios recién registrados, a
cambio de más falsos positivos. Los filtros de correo actúan antes de que el mensaje
llegue, pero como documenta Unit~42 \cite{unit42_2023qr} no leen el QR y quedan fuera de
juego. Las advertencias del navegador \cite{egelman2008warned,akhawe2013alice} actúan
después de navegar, cuando la petición al servidor del atacante ya salió. Y las apps
comerciales de escaneo seguro actúan en el punto correcto, entre el escaneo y la
apertura, pero con una detección insuficiente.

De esta comparación sale el hallazgo central del capítulo: 	extbf{el punto donde la
intervención es posible coincide con el punto donde ningún trabajo revisado aplica al
mismo tiempo detección multicapa y una comunicación bien fundamentada del veredicto.} El
momento entre leer el código y la redirección es el único en que la URL se puede ver, el
ataque todavía no se consuma y el usuario está en posición de decidir. Los sistemas de
reputación saben detectar pero no operan ahí; las apps comerciales operan ahí pero sin
detección suficiente ni criterio para comunicar; y la literatura de \emph{usable
security} sabe cómo comunicar pero no ha estudiado este contexto.

	extbf{Por cómo tratan la comunicación del veredicto} el desbalance es todavía más
marcado. De todos los antecedentes revisados, solo los de \emph{usable security}
\cite{cranor2008framework,egelman2008warned,akhawe2013alice} se ocupan explícitamente de
cómo presentar un resultado de seguridad, y ninguno lo hace en el contexto móvil de
escaneo de QR. Los trabajos técnicos de detección
\cite{sahoo2017malicious,sheng2009empirical,ma2009beyond,le2011phishdef} terminan en la
métrica de clasificación y no entran en la interfaz. Los proveedores de reputación
\cite{google_safebrowsing,virustotal_apiv3,phishtank_devinfo,urlhaus_api} entregan un
veredicto estructurado y dejan toda la presentación en manos de quien lo integre. Esa
separación entre quien detecta y quien comunica es justamente donde se ubica el aporte de
este trabajo.

La siguiente tabla resume la comparación de los principales antecedentes según su
enfoque, lo que aportan y lo que les falta para los fines de este proyecto.

\begin{longtable}{@{}>{\raggedright\arraybackslash}p{2.6cm}
                     >{\raggedright\arraybackslash}p{2.8cm}
                     >{\raggedright\arraybackslash}p{4.2cm}
                     >{\raggedright\arraybackslash}p{4.0cm}@{}}
\caption{Comparación de los trabajos revisados según enfoque, aporte y limitación.}
\label{tab:comparacion-antecedentes} \\
\toprule
\textbf{Trabajo / fuente} & \textbf{Enfoque metodológico} & \textbf{Aporte principal} & \textbf{Limitación relevante} \\
\midrule
\endfirsthead

\multicolumn{4}{@{}l}{\small\itshape (continuación de la Tabla \ref{tab:comparacion-antecedentes})} \\
\toprule
\textbf{Trabajo / fuente} & \textbf{Enfoque metodológico} & \textbf{Aporte principal} & \textbf{Limitación relevante} \\
\midrule
\endhead

\midrule
\multicolumn{4}{r@{}}{\small\itshape (continúa en la página siguiente)} \\
\endfoot

\bottomrule
\endlastfoot

Krombholz et al. \cite{krombholz2014qr} &
Taxonomía conceptual &
Formula el problema como cuestión de usabilidad; identifica la URL como vector dominante &
Anterior al escaneo nativo y a la adopción masiva post-2020 \\
\addlinespace

Vidas et al. \cite{vidas2013qrishing} &
Estudio de campo, conducta observada &
Evidencia cuantitativa de mayor susceptibilidad al QR frente a URL directa &
Muestra reducida; magnitudes no extrapolables al ecosistema actual \\
\addlinespace

Sahoo et al. \cite{sahoo2017malicious} &
Revisión sistemática &
Taxonomía de técnicas que fundamenta la arquitectura multicapa F1--F4 &
Tratamiento desactualizado del aprendizaje profundo \\
\addlinespace

Sheng et al. \cite{sheng2009empirical} &
Evaluación empírica de sistemas &
Establece el \emph{zero-hour phishing} como límite estructural de las listas &
Cifra del 20~\% no transferible a plataformas actuales \\
\addlinespace

Ma et al. \cite{ma2009beyond} &
Clasificación sobre conjunto de datos &
Precisión cercana al 99~\% sin visitar el sitio destino &
Condiciones de laboratorio; sin despliegue móvil \\
\addlinespace

Le et al. \cite{le2011phishdef} &
Clasificación léxica en línea &
Precisión superior al 95~\% solo con características de la URL &
Vulnerable a acortadores y a dominios comprometidos \\
\addlinespace

Google Safe Browsing \cite{google_safebrowsing} &
Documentación técnica oficial &
Mayor cobertura indexada; latencia inferior a 200~ms &
Criterio de clasificación opaco, no auditable \\
\addlinespace

VirusTotal \cite{virustotal_apiv3} &
Documentación técnica oficial &
Agregación de más de 70 motores; reduce falsos negativos &
Plan gratuito: 4 consultas/min y 500/día \\
\addlinespace

PhishTank \cite{phishtank_devinfo} / URLhaus \cite{urlhaus_api} &
Plataforma colaborativa &
Rapidez ante campañas emergentes &
Cobertura desigual por región e idioma \\
\addlinespace

Cranor \cite{cranor2008framework} &
Modelo conceptual &
Marco diagnóstico de los tres eslabones del comportamiento seguro &
No operacionalizado como instrumento de medición \\
\addlinespace

Egelman et al. \cite{egelman2008warned} &
Experimento controlado &
Criterios de diseño de advertencias basados en evidencia &
Contexto exclusivamente de navegador de escritorio \\
\addlinespace

Akhawe y Felt \cite{akhawe2013alice} &
Estudio de campo a gran escala &
La severidad percibida modula la respuesta: tasa de continuación del 70~\% en SSL frente al 25~\% en malware &
No validado en contexto móvil \\
\addlinespace

ISO/IEC 18004 \cite{iso18004_2015,iso18004_2024} &
Estándar internacional &
Delimita el contenido codificable y la lógica de clasificación &
No define autenticación del emisor ni integridad del contenido \\
\addlinespace

OWASP MASVS \cite{owasp2024masvs} &
Estándar de verificación &
Criterios verificables de seguridad de la aplicación &
Estándar de requisitos, no metodología de implementación \\
\addlinespace

\emph{Safe QR Scan}, \emph{SafeQR Reader} (App Store) &
Producto de mercado; sin documentación técnica pública &
Validan la viabilidad y la aceptación del punto de intervención &
Cerradas e inauditables; detección basada en lista negra \\

\end{longtable}

\section{Aportes y limitaciones de los trabajos revisados}

\subsection{Aportes principales}

Los antecedentes revisados le dan a este proyecto cuatro bases distintas. Del bloque
que caracteriza el vector \cite{krombholz2014qr,vidas2013qrishing} se toma la
confirmación de que el problema es de usabilidad y no solo técnico, junto con la
evidencia de que caer en \emph{quishing} es algo estructural y no un descuido
individual. Del bloque de detección
\cite{sahoo2017malicious,sheng2009empirical,ma2009beyond,le2011phishdef} se toma la
clasificación de técnicas que sustenta la arquitectura multicapa F1--F4, y la evidencia
de que ninguna técnica sola alcanza. Del bloque de infraestructura
\cite{google_safebrowsing,virustotal_apiv3,phishtank_devinfo,urlhaus_api,misp_project}
se toma que existen fuentes de reputación consultables en tiempo real con latencias
compatibles con el objetivo de cinco segundos, además del modelo colaborativo en que se
apoya la función de reporte del prototipo. Del bloque de \emph{usable security}
\cite{cranor2008framework,egelman2008warned,akhawe2013alice} se toman los criterios de
diseño de advertencias que sustentan mostrar el veredicto en dos niveles de detalle. Y
los estándares \cite{owasp2024masvs,iso18004_2015,iso18004_2024} junto con la
documentación técnica
\cite{apple_vndetectbarcodes,google_mlkit_barcode,rousavy_visioncamera,apple2024platformsecurity,apple_ats,android_netsecconfig}
aportan los criterios verificables de cumplimiento y la confirmación de que procesar
todo en el teléfono es viable.

\subsection{Limitaciones identificadas}

Mirando críticamente esos mismos trabajos aparecen seis limitaciones que, juntas,
delimitan el vacío de investigación.

\paragraph{Desactualización de la evidencia empírica sobre el vector.}
Los estudios de referencia sobre susceptibilidad
\cite{krombholz2014qr,vidas2013qrishing} son anteriores a que el escaneo de QR se
integrara a la cámara del teléfono y a la adopción masiva posterior a 2020. Lo que
concluyen sigue siendo válido en lo cualitativo, pero sus números ya no se pueden
trasladar al escenario actual.

\paragraph{Transferencia no validada de la evidencia sobre advertencias.}
Los hallazgos sobre qué tan efectivas son las advertencias
\cite{egelman2008warned,akhawe2013alice} vienen todos de navegadores de escritorio.
Aplicarlos al escaneo móvil (pantalla chica, uso en movimiento, sin barra de
direcciones) es una extrapolación razonable que ninguna publicación revisada ha
comprobado.

\paragraph{Brecha entre desempeño de laboratorio y despliegue real.}
Los clasificadores léxicos reportan precisiones del 95~\% al 99~\%
\cite{ma2009beyond,le2011phishdef} sobre conjuntos de datos históricos, sin evidencia de
cómo se comportan bajo las restricciones de cómputo, batería y conectividad de un
teléfono en uso real.

\paragraph{Ausencia de evaluación independiente de las fuentes de reputación.}
No se encontró literatura arbitrada que compare qué tanto coinciden Google Safe
Browsing, PhishTank y VirusTotal sobre un mismo conjunto de URLs, ni que describa cómo se
comporta un veredicto combinado. Lo único disponible es la documentación de cada
proveedor.

\paragraph{Opacidad de las soluciones comerciales.}
Las apps existentes no publican su arquitectura, ni sus fuentes, ni sus tasas de
error, así que no se pueden auditar. Eso impide compararlas científicamente y también
aprender de las decisiones que tomaron.

\paragraph{Fragmentación entre detección y comunicación.}
Ningún trabajo revisado junta en una misma propuesta la detección multicapa y un
diseño fundamentado de cómo comunicar el veredicto. La literatura técnica termina en la
métrica; la de usabilidad arranca cuando la advertencia ya está formulada. El espacio de
en medio, o sea convertir un veredicto técnico en una decisión que un usuario común pueda
tomar, en el teléfono y en menos de cinco segundos, sigue sin abordarse de forma
integrada.

A estas se suma una limitación de contexto: 	extbf{no hay producción académica sobre
el fenómeno en el entorno latinoamericano y guatemalteco}. Toda la literatura revisada
viene de contextos anglosajones y europeos, mientras que las plataformas colaborativas de
reporte seguramente cubren menos las campañas en español dirigidas a la región. Eso
sugiere que el límite de las listas negras podría ser incluso más grave aquí que donde se
midió.

\section{Fundamentación del proyecto propuesto}

Los antecedentes sustentan el proyecto en cuatro planos, y conviene dejar clara la
correspondencia entre cada limitación identificada y la decisión de diseño que la
responde.

\paragraph{Fundamentación del punto de intervención.}
Cuatro puntos de evidencia apuntan a lo mismo. Krombholz et al.
\cite{krombholz2014qr} muestran que no se puede inspeccionar el contenido antes de
escanear. Vidas et al. \cite{vidas2013qrishing} miden qué tan susceptible resulta la
gente por eso. Unit~42 \cite{unit42_2023qr} documenta que los filtros de correo no ven el
QR. Y la falta de mecanismos de autenticación en ISO/IEC~18004
\cite{iso18004_2015,iso18004_2024} impide resolverlo en la capa del código. Todo junto
indica que el único punto de la cadena del ataque que sigue sin protección automatizada
es el teléfono del usuario en el instante del escaneo. El objetivo general del trabajo,
meterse entre la lectura del código y la redirección, sale directamente de esa evidencia
y no de una preferencia de diseño.

\paragraph{Fundamentación de la arquitectura multicapa.}
La clasificación de Sahoo et al. \cite{sahoo2017malicious} y la evidencia sobre los
límites de cobertura de las listas \cite{sheng2009empirical}, puesta en contexto con los
datos de APWG de 2024 y 2025 \cite{apwg2024q1,apwg2025q1}, sustentan que ninguna técnica
sola resuelve el problema. De ahí la arquitectura F1--F4. Que F1 sea el primer filtro se
apoya en lo que ya se sabe de las fuentes de reputación
\cite{google_safebrowsing,virustotal_apiv3,phishtank_devinfo}: responden rápido, aciertan
mucho sobre amenazas conocidas y cuestan poco en cómputo, lo que las vuelve el punto de
decisión inicial más lógico bajo el límite de cinco segundos. Y saber que no alcanzan es,
al mismo tiempo, la razón de que haya filtros después.

\paragraph{Fundamentación de las decisiones de interfaz.}
El marco de Cranor \cite{cranor2008framework} y los resultados de Egelman et al.
\cite{egelman2008warned} y de Akhawe y Felt \cite{akhawe2013alice} sustentan mostrar el
veredicto en dos niveles de detalle con divulgación progresiva, usar lenguaje concreto en
vez de técnico, diferenciar visualmente los niveles de riesgo e interrumpir el flujo a
propósito. Que la tasa de continuación varíe entre el 25~\% y el 70~\% según el tipo de
amenaza \cite{akhawe2013alice} sustenta específicamente que los tres veredictos (seguro,
sospechoso, malicioso) tengan que distinguirse sin ambigüedad y no ser matices de una
misma escala.

\paragraph{Fundamentación de las decisiones técnicas.}
La documentación de las APIs nativas de visión
\cite{apple_vndetectbarcodes,google_mlkit_barcode} y de react-native-vision-camera
\cite{rousavy_visioncamera} confirma que procesar en el teléfono con menos de 100~ms es
viable, condición necesaria para el presupuesto de tiempo total. El MASVS
\cite{owasp2024masvs} sustenta los controles de almacenamiento cifrado, seguridad de red
y permisos mínimos, bajo un criterio verificable: una app que existe para proteger al
usuario tiene que ser ella misma auditablemente segura.

\paragraph{Aporte diferencial.}
Con todo esto sobre la mesa, el aporte de este trabajo se puede definir con precisión
y sin inflarlo. No consiste en proponer una técnica nueva de detección, porque el
proyecto integra fuentes que ya existen, ni en superar a las apps comerciales,
comparación que quedó fuera del alcance. Consiste en tres cosas: (i)~integrar de forma
documentada y auditable, en el punto correcto, la detección multicapa y una comunicación
fundamentada del veredicto, lo que responde a la sexta limitación; (ii)~aportar evidencia
empírica sobre cómo se comporta un veredicto combinado de tres fuentes de reputación, lo
que responde a la cuarta; y (iii)~aplicar los principios de \emph{usable security} al
escaneo móvil de QR, documentando la extrapolación que la segunda limitación deja
pendiente. Las restricciones también son explícitas: solo F1 se implementa de extremo a
extremo, la validación de usabilidad no alcanza representatividad estadística y no se
compara con productos comerciales.

\section{Conclusión: vacíos identificados y necesidad del proyecto}

Revisar críticamente los antecedentes lleva a una conclusión: hoy el \emph{quishing}
no está limitado por falta de capacidad técnica para detectar. Esa capacidad existe,
funciona a escala mundial y se puede consultar por interfaces públicas en milisegundos.
Lo que falta es aplicarla en el lugar y en el momento donde el ataque se concreta, y
cerrar la distancia entre quienes desarrollan la detección y quienes estudian cómo
explicarle el resultado a un usuario que no es experto.

Los vacíos identificados son cuatro. Hay un 	extbf{vacío de aplicación}: la
infraestructura de reputación está desplegada en navegadores de escritorio y en filtros
de correo, pero no en el escaneo móvil, que es justamente el único punto donde la URL de
un QR se vuelve legible antes de navegar. Hay un 	extbf{vacío de integración}: la
literatura técnica de detección y la de \emph{usable security} avanzaron sin cruzarse, y
por eso ninguna propuesta revisada aborda de forma unificada la detección multicapa junto
con el diseño de cómo comunicar el veredicto. Hay un 	extbf{vacío de evaluación}: no hay
evidencia arbitrada sobre qué tanto coinciden las principales fuentes de reputación ni
sobre cómo se comporta un veredicto combinado, ni tampoco validación de los hallazgos
sobre advertencias en contexto móvil. Y hay un 	extbf{vacío de transparencia}: las
únicas soluciones que operan en el punto correcto son productos cerrados, imposibles de
auditar y con un mecanismo insuficiente, lo que impide que el conocimiento sobre este
problema se acumule.

Estos vacíos justifican el proyecto. Desarrollar un prototipo de app móvil que se
meta antes de la redirección al momento de escanear, que combine en tiempo real los
veredictos de tres fuentes independientes de reputación con un límite de cinco segundos,
y que muestre el resultado en una interfaz diseñada según los principios documentados de
\emph{usable security}, responde directamente a cada uno de ellos. Documentar de forma
abierta su arquitectura, sus fuentes y sus resultados de validación, incluidas las
limitaciones, atiende además el vacío de transparencia y deja la puerta abierta para que
otros trabajos extiendan la propuesta hacia los filtros heurístico, de \emph{sandbox} y de
aislamiento de navegación que ya están previstos en el diseño.

El contexto refuerza la pertinencia del esfuerzo. Toda la literatura revisada viene de
contextos anglosajones y europeos, mientras que la evidencia disponible sugiere que en
Latinoamérica los códigos QR se adoptaron más rápido de lo que crecieron las prácticas de
seguridad, y que las plataformas colaborativas de reporte cubren menos las campañas
dirigidas a la región. Un trabajo que documente el problema y proponga una solución
evaluable desde Guatemala no solo llena un vacío técnico, sino que deja un antecedente en
un área donde la producción académica local todavía es escasa.